% !PRINTED 195
\weberpage{195}
\begin{center}
\S.\,51.

\gesperrt{Finite groups of linear fractional substitutions.}

\gesperrt{Poles of the groups.}
\end{center}

From the developments so far it follows that, once all finite
groups of linear fractional substitutions have been found, all
finite groups of properly orthogonal ternary substitutions, and
all finite groups of binary linear substitutions with determinant
1, can be found from them without further ado.

Besides these there are also finite groups which contain,
along with the properly orthogonal substitutions, improperly
orthogonal ones as well.

These groups, to which we shall return later, cannot be
derived from the groups of linear fractional substitutions alone,
because with the fractional substitutions the distinction between
properly and improperly orthogonal substitutions is effaced.

We therefore now seek first of all to determine \gesperrt{all finite
groups of linear fractional substitutions} $^{1)}$.

We denote by $G$ such a group of degree $n$,
and by
\weq{1}{x,\ \Theta_1(x),\ \Theta_2(x),\ \ldots,\ \Theta_{n-1}(x)}
its elements, in which the $\Theta$ are symbols for linear functions
\weq{2}{\Theta(x) = \frac{\alpha x + \beta}{\gamma x + \delta}}
for which the notation
\weq{3}{\Theta = \begin{pmatrix}\alpha, & \beta\\ \gamma, & \delta\end{pmatrix}}
is also used.

From every group of the form (1) we can, by \S.\,37, derive
a whole family of isomorphic groups, if we denote by
$L$ an arbitrary linear substitution:
\weq{4}{1,\ L\Theta_1L^{-1},\ L\Theta_2L^{-1},\ \ldots,\ L\Theta_{n-1}L^{-1},}
%CONT>>

\par\noindent\rule{4cm}{0.4pt}\par
$^{1)}$ On the theory of these groups compare especially:
\gesperrt{Gordan}, Ueber endliche Gruppen linearer Transformationen einer
Veränderlichen. Mathematische Annalen, Bd. XII, S. 23 (1877);
\gesperrt{Klein}, Vorlesungen über das Ikosaëder (Leipzig 1884).
